% -------------------------------------- Modèles : Partie électrochimique -----------------------------------------

\begin{frame}[t]{Modèles : partie électrochimique}
  \vspace{-1em}

  \begin{columns}[T,onlytextwidth]

    % ----------- Schéma -----------
    \column{0.48\textwidth}
      \begin{pccard}
        \centering
        \figschemaelectrochimique
        \vspace{0.5em}
        {\scriptsize\bfseries\color{pcNavy}
        Représentation de la partie électrochimique}
      \end{pccard}

    % ----------- Différents angles d'approche -----------
    \column{0.48\textwidth}
      \vspace{-0.5em}
      \begin{pcbanner}{Différents angles d'approche}

        %\vspace{-0.25em}

        \begin{pcbannerlight}{Apport du CO2}
          \vspace{0.3em}
          \begin{itemize}\footnotesize
                \item Bullage dans l'électrolyte 
                \item Electrode à diffusion de gaz (GDE)
          \end{itemize}
        \end{pcbannerlight}

        \vspace{-0.5em}

        \begin{pcbannerlight}{Géométrie de l'électrode}
          \vspace{0.3em}
          \begin{itemize}\footnotesize
            \item Modèle 3D local 
            \item Modèle 1D macroscopique
          \end{itemize}
        \end{pcbannerlight}

      \end{pcbanner}
    \end{columns}


    \vspace{-0.75em}
    \begin{columns}[T,onlytextwidth]
      \column{\textwidth}
        \begin{pcbanner}{Démarche}
          \begin{itemize}\footnotesize
            \item Comparer 1D et 3D (homogénéisation)
            \item Comparer les deux configurations (électrode à diffusion de gaz / bullage dans l'électrolyte)
            \item Confronter à l'expérimental (partenaires GASPER)
          \end{itemize}
        \end{pcbanner}
    \end{columns}

\end{frame}


%--------------------------------------------------------------------
\begin{frame}[t]{Modèles : partie électrochimique}
  \vspace{-1em}

  \begin{columns}[T,onlytextwidth]

    % ----------- Schéma (inchangé) -----------
     \column{0.48\textwidth}
      \begin{pccard}
        \centering
        \figschemaelectrochimique
        \begin{tikzpicture}[remember picture, overlay]
          \draw[pcSun, line width=1.5pt]
            ([xshift=-4.35cm, yshift=0.91cm]current page.center)
            ellipse [x radius=0.6cm, y radius=1.1cm];
        \end{tikzpicture}
        \vspace{0.5em}
        {\scriptsize\bfseries\color{pcNavy}
        Représentation de la partie électrochimique}
      \end{pccard}

      \begin{pcbanner}{Réduction du \ch{CO2}}
      \centering\scriptsize
      \textbf{Équation de réaction}
      \[ \mathrm{CO_2} + 2\,\mathrm{H^+} + 2\,\mathrm{e^-} \rightarrow \mathrm{CO} + \mathrm{H_2O} \]
      \end{pcbanner}

    % ----------- Modèle électrochimique -----------
    \column{0.48\textwidth}
      \vspace{-0.5em}
        \begin{pcbanner}{Modèle dans le milieu poreux}
        \vspace{0.4em}
        \centering\scriptsize
        \textbf{Équation de conservation}
        \[ \frac{\partial C_i}{\partial t} + \nabla\!\cdot\! N_i = 0 \]
        \textbf{Transport des espèces (Nernst--Planck)}
        \[ N_i = -D_i\nabla C_i - \frac{z_i F}{RT}\,D_i C_i\,\nabla\varphi_l + C_i\,\vec{v} \]

        \textbf{Électroneutralité} \qquad \textbf{Transport électronique (Ohm)}
        \[ \sum_i z_i C_i = 0 \qquad \qquad \qquad \nabla\!\cdot\!\left(-\sigma_s\,\nabla\varphi_s\right) = 0 \]

        \textbf{Cinétique interfaciale}\\[0.3em]
        \begin{itemize}\scriptsize
          \item Butler--Volmer (semi-empirique)
          \item Marcus--Gerischer (théorique)
        \end{itemize}
        \end{pcbanner}

    \end{columns}

\end{frame}

%--------------------------------------------------------------------
\begin{frame}[t]{Electrochimie : cas simplifié \ch{CO2} -> \ch{CO}}
  \vspace{-1em}

  \begin{columns}[T,onlytextwidth]

    % ----------- Colonne gauche : hypothèses -----------
      \column{0.48\textwidth}
\vspace{-0.5em}

   \begin{pcbanner}{Module de Thiele $\phi$}
        \footnotesize
        Compétition diffusion / réaction :
        \begin{itemize}\footnotesize
        \item $\phi \ll 1$ : diffusion rapide $\rightarrow$ concentration uniforme
        \item $\phi \gg 1$ : réaction rapide $\rightarrow$ \ch{CO2} consommé près de l'interface
        \end{itemize}
      \end{pcbanner}

      \begin{pcbanner}{Hypothèses (approche de Thiele)}
        \vspace{0.15em}
        \begin{itemize}\footnotesize
          \item Régime stationnaire, transport purement diffusif
          \item Migration, convection, chute ohmique négligées
          \item Deux espèces carbonées seulement : \ch{CO2}, \ch{CO}
          \item Réaction globale simplifiée \ch{CO2 -> CO}
        \end{itemize}
      \end{pcbanner}

    % ----------- Colonne droite : modèle + solution -----------
    \column{0.48\textwidth}
      \vspace{-0.5em}
      \begin{pcbanner}{Système diffusion--réaction}
        \centering\scriptsize
       \textbf{Forme adimensionnée (module de Thiele $\phi$)}
        \[ \frac{\mathrm{d}^2 \langle C_{\ch{CO2}}\rangle^l}{\mathrm{d}\zeta^2}
           = \phi^2\, \langle C_{\ch{CO2}}\rangle^l \]
        \[ \frac{\mathrm{d}^2 \langle C_{\ch{CO}}\rangle^l}{\mathrm{d}\zeta^2}
           = -\,\frac{a_v^{\mathrm{cat}}\, k\, l_e^2}{\varepsilon_l D^{\mathrm{eff}}_{\ch{CO}}}
             \, \langle C_{\ch{CO2}}\rangle^l \]
      \end{pcbanner}

      \begin{pcbanner}{Solution analytique}
        \vspace{-0.5em}
        \centering\scriptsize
        \[ \langle C_{\ch{CO2}}\rangle_l(\zeta)
           = C^{\mathrm{hom}}_{\ch{CO2}}\,
             \frac{\cosh[\phi(\zeta+1)]}{\cosh(\phi)} \]
        \[ \langle C_{\ch{CO}}\rangle_l(\zeta)
           = C^{\mathrm{hom}}_{\ch{CO}}
           + \frac{D^{\mathrm{eff}}_{\ch{CO2}}}{D^{\mathrm{eff}}_{\ch{CO}}}
             \,C^{\mathrm{hom}}_{\ch{CO2}}
             \left[1 - \frac{\cosh[\phi(\zeta+1)]}{\cosh(\phi)}\right] \]
      \end{pcbanner}

    \end{columns}

\end{frame}